\documentclass[aspectratio=169, 10pt]{beamer}
\usetheme{metropolis}

% --- Edinburgh colours ----------------------------------------
\definecolor{UoERed}{HTML}{7A2318}
\definecolor{UoEGold}{HTML}{B8860B}
\definecolor{CodeBG}{HTML}{F7F3ED}
\definecolor{CodeFrame}{HTML}{D5C9B1}

\setbeamercolor{palette primary}{bg=UoERed, fg=white}
\setbeamercolor{frametitle}{bg=UoERed, fg=white}
\setbeamercolor{title separator}{fg=UoEGold}
\setbeamercolor{progress bar}{fg=UoEGold, bg=UoEGold!20}
\setbeamercolor{alerted text}{fg=UoERed}
\setbeamercolor{block title}{bg=UoERed!15, fg=UoERed}
\setbeamercolor{block body}{bg=UoERed!5}

% --- Packages -------------------------------------------------
\usepackage[T1]{fontenc}
\usepackage{lmodern}
\usepackage{amsmath, amssymb, mathtools}
\usepackage{listings}
\usepackage{graphicx, booktabs}
\usepackage{tikz}
\usetikzlibrary{arrows.meta, positioning, calc, decorations.pathreplacing}
\usepackage{hyperref}
\hypersetup{colorlinks=true, linkcolor=UoERed, urlcolor=UoERed!70!black}
\setbeamertemplate{navigation symbols}{}

\lstset{
  language=Python,
  basicstyle=\ttfamily\scriptsize,
  keywordstyle=\color{UoERed}\bfseries,
  commentstyle=\color{gray}\itshape,
  stringstyle=\color{UoEGold},
  backgroundcolor=\color{CodeBG},
  frame=single,
  rulecolor=\color{CodeFrame},
  numbers=none, breaklines=true, showstringspaces=false, tabsize=4,
  xleftmargin=4pt, xrightmargin=4pt, aboveskip=6pt, belowskip=4pt,
}

\AtBeginSection[]{
  \begin{frame}
    \vfill\centering
    \usebeamerfont{title}\insertsectionhead\par
    \vfill
  \end{frame}
}

\title{Week 1 --- From Numerical Methods to Neural Networks}
\subtitle{Why Traditional Methods Struggle \& Neural Networks as Function Approximators}
\author{Dr Juan Zurita}
\institute{Deep Learning for Macroeconomics\\
The University of Edinburgh~$\cdot$~School of Economics}
\date{}

\begin{document}
\maketitle

\begin{frame}{The Challenge}
Traditional numerical methods break down in \alert{high dimensions}.

\bigskip
\begin{block}{The Curse of Dimensionality}
A grid with $n$ points per dimension in $d$ dimensions has $n^d$ nodes.

$n=100$, $d=5$ $\Rightarrow$ $10^{10}$ grid points --- infeasible.
\end{block}

\bigskip
Deep learning offers a \textbf{grid-free} alternative.
\end{frame}

\begin{frame}{From PNM to Deep Learning}
\begin{center}
\begin{tabular}{ll}
\toprule
PNM & This course \\
\midrule
Root-finding & Equilibrium as loss functions \\
BFGS / Nelder--Mead & Gradient descent, Adam \\
Polynomials / splines & Neural networks \\
Value-function iteration & Deep Equilibrium Networks \\
Grid-based methods & Grid-free solutions \\
\bottomrule
\end{tabular}
\end{center}

\bigskip
The key shift: parameterise value and policy functions as \alert{neural networks}.
\end{frame}

\begin{frame}{Universal Approximation Theorem}
\begin{block}{Cybenko (1989), Hornik et al.\ (1989)}
A feedforward network with one hidden layer and a non-linear activation can approximate any continuous function on a compact set to arbitrary accuracy.
\end{block}

\bigskip
This is the theoretical foundation for using neural networks as function approximators in economics.

\bigskip
Practical question: \textbf{how many} neurons do we need?
\end{frame}

\begin{frame}{A Single Neuron}
$$y = \sigma\!\left(\sum_{i=1}^n w_i x_i + b\right)$$

\bigskip
\begin{itemize}
\item $\mathbf{w}$: weights (parameters to learn)
\item $b$: bias
\item $\sigma$: activation function (ReLU, sigmoid, tanh)
\end{itemize}

\bigskip
A \textbf{layer} stacks many neurons; a \textbf{network} stacks many layers.
\end{frame}

\begin{frame}{Your First Network in PyTorch}
\begin{block}{Key idea}
Define the architecture, choose a loss, run gradient descent.
\end{block}

\bigskip
In the notebook you will:
\begin{enumerate}
\item Approximate $\sin(x)$ with a small network
\item Compare with polynomial and Chebyshev approximation from PNM
\item See how adding layers and neurons improves accuracy
\end{enumerate}

\bigskip
\textbf{Next week:} the full deep learning toolkit --- backpropagation, regularisation, and training.
\end{frame}

\end{document}
